\documentclass[a4paper,12pt]{article}
\usepackage[utf8]{inputenc}
\usepackage[slovak]{babel}
\usepackage{hyperref}
\usepackage{graphicx}
\usepackage{fancyhdr}
\usepackage{titlesec}
\usepackage{amsmath}
\usepackage[utf8]{inputenc}

% Nastavenie titulnej strany
\begin{document}

\begin{titlepage}
    \centering
    \vspace*{1cm}
    \Large{\textbf{Technická univerzita v Košiciach}\\
    Fakulta elektrotechniky a informatiky}\\
    \vfill
    \large{Formálne jazyky}\\
    \large{Dokumentácia k obhajobe zadania}
    \vfill
    \begin{flushleft}
        Dmytro Shkyl\\
        Akademický rok: 2024/2025\\
        Študijný program: informatika
    \end{flushleft}
\end{titlepage}

\tableofcontents
% Nastavenie formulacie zadania
\newpage
\section{Formulácia zadania}

\indent\textbf{Téma:} Generátor konečnostavových automatov\\

\noindent\textbf{Cieľ zadania:} Implementovať generátor nedeterministických konečnostavových automatov (NKA) na základe zadaného regulárneho výrazu. Pri implementácii je potrebné využiť metódu syntaktickej analýzy zhora nadol (rekurzívny zostup).

\vspace{1em}
\textbf{Požiadavky na vstup:}
\begin{itemize}
    \item Vstupom je reťazec znakov, ktorý reprezentuje regulárny výraz v štandardne využívanej notácii pre predmet Formálne jazyky.
\end{itemize}

\vspace{1em}
\textbf{Výstup programu:}
\begin{itemize}
    \item Grafická reprezentácia odvodeného stromu pre zadaný regulárny výraz. 
    \item Plne funkčná programová implementácia nedeterministického konečnostavového automatu zodpovedajúceho zadanému regulárnemu výrazu.
\end{itemize}
% Vypracovanie zadania
\newpage
\section{Vypracovanie zadania}

\subsection{Opis architektúry riešenia}
Architektúra je navrhnutá modulárne, aby bolo možné  jednotlivé časti programu samostatne testovať.
Navrhnuté riešenie je rozdelené do štyroch hlavných modulov:

\begin{itemize}
    \item \textbf{Lexer} \\
    Lexer zodpovedá za rozdelenie vstupného reťazca (regulárneho výrazu) na tokeny. Spracováva jednotlivé symboly.
    \item \textbf{Parser} \\
    Parser prijíma tokeny z lexera a pomocou metódy rekurzívneho zostupu vytvára syntaktický strom, ktorý reprezentuje štruktúru regulárneho výrazu.

    \item \textbf{Generátor NKA} \\
    Tento modul vytvára nedeterministický konečnostavový automat (NKA) na základe syntaktického stromu. Každá operácia regulárneho výrazu má definované pravidlá.

    \item \textbf{Main} \\
    Modul \texttt{main} zabezpečuje spustenie programu, načítanie vstupného výrazu, inicializáciu lexera, parsera a generátora NKA. Tiež môže zobrazovať grafickú reprezentáciu odvodeného stromu.
\end{itemize}

\bigskip



\subsection{Gramatika jazyka regulárnych výrazov}

Gramatika pre regulárne výrazy môže byť definovaná ako:

\begin{itemize}
    \item \textbf{Symboly:} 
    \begin{equation*}
        R \rightarrow a \mid b \mid c \mid \dots \quad (\text{symbole ako } a, b, c, \dots)
    \end{equation*}
    \item \textbf{Postupnost':} 
    \begin{equation*}
        R \rightarrow R R \quad (\text{sekvencia dvoch symbolov})
    \end{equation*}
    \item \textbf{Alternatíva:} 
    \begin{equation*}
        R \rightarrow R \mid R \quad 
    \end{equation*}
    \item \textbf{Pozitívny uzáver:} 
    \begin{equation*}
        R \rightarrow R\{R\} \quad (\text{jeden alebo viac opakovaní}) 
    \end{equation*}
    \item \textbf{Tranzitívny uzáver:} 
    \begin{equation*}
        R \rightarrow \{R\} \quad (\text{opakovanie nula alebo viackrát}) 
    \end{equation*}
    \item \textbf{Voliteľnosť:}
    \begin{equation*}
        R \rightarrow [R] \quad (\text{nula alebo jedenkrát})
    \end{equation*}
\end{itemize}

\subsection{Podrobný postup implementácie}
1. \textbf{Lexikálny analyzátor (Lexer)}

Lexikálny analyzátor je zodpovedný za rozdelenie vstupného reťazca na jednotlivé tokeny, ktoré sú základom pre ďalší syntaktický rozbor. Implementácia lexikálneho analyzátora zahŕňa:

\begin{itemize}
    \item \textbf{Inicjalizácia}: Lexer sa inicializuje s reťazcom, ktorý sa má analyzovať. Určuje sa počiatočná pozícia na začiatok reťazca.
    \item \textbf{Metóda \texttt{advance()}}: Posúva aktuálnu pozíciu o jeden znak v reťazci a aktualizuje aktuálny znak. Ak sa pozícia dostane za koniec reťazca, nastaví sa aktuálny znak na \texttt{None}.
    \item \textbf{Metóda \texttt{skip\_whitespace()}}: Preskakuje všetky medzery a nové riadky, aby analyzátor mohol pokračovať v spracovaní tokenov.
    \item \textbf{Metóda \texttt{get\_next\_token()}}: Hlavná metóda lexikálneho analyzátora, ktorá prechádza vstupný reťazec a vracia nasledujúci token podľa rozpoznaných symbolov.
\end{itemize}

Výsledkom lexikálneho analyzátora sú tokeny, ako napríklad symboly (písmená), operátory, skratky pre zátvorky a iné špeciálne znaky.

2. \textbf{Syntaktický analyzátor (Parser)}

Syntaktický analyzátor spracováva sekvenciu tokenov v syntaktické stromy, ktoré reprezentujú štruktúru pravidiel pre dané regulárne výrazy. Implementácia je založená na algoritme rekursívneho spúšťania, ktorý sa skladá z niekoľkých metód pre analýzu rôznych častí regulárnych výrazov.

\begin{itemize}
     \item \textbf{Metóda \texttt{Alternative()}}: Táto metóda analyzuje alternatívy (výrazy spojené operátorom \texttt{|}), ale ak sa v rámci regulárneho výrazu nenájde operátor \texttt{|}, vráti sekvenciu. Rekurzívne volá metódu \texttt{Sequence()} na analyzovanie jednotlivých alternatív. Ak sa zistí \texttt{|}, vytvorí sa nový uzol v strome pre alternatívu.
    \item \textbf{Metóda \texttt{Sequence()}}: Analyzuje sekvencie, kde výrazy musia byť spracované v presnom poradí. Umožňuje kombinovať viaceré faktory do sekvencie.
    \item \textbf{Metóda \texttt{Factor()}}: Analyzuje jednotlivé faktory, ako sú symboly, zátvorky. Na základe tokenu rozhodne, ako tento token spracovať.
\end{itemize}

Pri spracovaní regulárnych výrazov sa používajú rôzne typy uzlov, ktoré reprezentujú elementy regulárnych výrazov. Tieto uzly sú:
\begin{itemize}
    \item \textbf{regular}: reprezentuje regulárny výraz.
    \item \textbf{element}: reprezentuje symboly.
    \item \textbf{alternative}: reprezentuje alternatívu (výrazy oddelené operátorom \texttt{|}).
    \item \textbf{sequence}: reprezentuje sekvenciu (výrazy vykonané v poradí).
    \item \textbf{LCBRA a RCBRA}:  reprezentuje opakovanie (výraz v uzavretych zátvorkách \texttt{\{\}}).
    \item \textbf{LBRCKT a RBRCKT}: reprezentuje volitelnost' (výraz v hranatých zátvorkách \texttt{[]}).
\end{itemize}

3. \textbf{Vytváranie a vizualizácia stromu}

Na vizualizáciu syntaktického stromu sa použila knižnica \texttt{anytree}. Vizualizácia umožňuje zobrazovať štruktúru regulárneho výrazu v podobe stromu s uzlami reprezentujúcimi jednotlivé časti výrazu. Tento proces zahŕňa:

\begin{itemize}
    \item \textbf{Metóda \texttt{has\_pipe()}}: Prechádza strom a zisťuje, či existuje operátor \texttt{|}. To je dôležité pre správne vykreslenie alternatívnych výrazu.
    \item \textbf{Metóda \texttt{build\_tree()}}: Rekurzívne konštruuje strom, pričom pre každý uzol vyhodnocuje jeho typ a pridáva príslušné deti do stromu.
\end{itemize}

V prípade, že strom obsahuje operátor \texttt{|}, vytvorí sa nový uzol typu \texttt{alternative}, ktorý je potom napojený na ďalšie uzly.


\subsection{Ukážka vstupu a prezentácia výsledkov}
Pre spracovanie regulárnych výrazov sa používa rôzne výrazy, ktoré sú spracované na syntaktické stromy. Príklad vstupu:

\newpage
\begin{itemize}
    \item \texttt{\{ab|c\}}: Tento regulárny výraz sa používa na opakovanie reťazcov, ktoré začínajú písmenom \texttt{a}, následne obsahujú \texttt{b} alebo \texttt{c}. Prazdny retazec tiez bude akceptovany.
\end{itemize}

\begin{verbatim}
regular
- alternative
    - sequence
        - element
            - <LCBRA>
            - regular
                - alternative
                    - sequence
                        - element
                            - SYMBOL
                                - a
                        - element
                            - SYMBOL
                                - b
                    - PIPE
                    - sequence
                        - element
                            - SYMBOL
                                - c
            - <RCBRA>
    \end{verbatim}

Pre spracovanie regulárneho výrazu \texttt{\{ab|c\}} môžeme generovať konečný automat (NKA), ktorý akceptuje nasledovné reťazce:

\begin{enumerate}
    \item \texttt{ab}: Akceptované
    \item \texttt{c}: Akceptované
    \item \texttt{a}: Neakceptované
    \item \texttt{abc}: Akceptované
    \item \texttt{} (prázdny reťazec): Akceptované
\end{enumerate}

\newpage
\begin{itemize}
    \item \texttt{[a][b][c]}: Tento regulárny výraz znamená, že každe pismeno môže prist' nula alebo jeden krat. Prazdny retazec tiez bude akceptovany.
\end{itemize}

\begin{verbatim}
  regular
     - sequence
     | - element
     |        - <LCBRCKT>
     |        - regular
     |             - element
     |                   - SYMBOL
     |                        - a
     |        - <RCBRCKT>
     |  - element
     |        - <LCBRCKT>
     |        - regular
     |             - element
     |                   - SYMBOL
     |                        - b
     |        - <RCBRCKT>
     |	- element
     |        - <LCBRCKT>
     |        - regular
     |             - element
     |                   - SYMBOL
     |                        - c
     |       - <RCBRCKT>
    \end{verbatim}

Pre spracovanie regulárneho výrazu \texttt{[a][b][c]} môžeme generovať konečný automat (NKA), ktorý akceptuje nasledovné reťazce:

\begin{enumerate}
    \item \texttt{ab}: Akceptované
    \item \texttt{c}: Akceptované
    \item \texttt{abc}: Akceptované
    \item \texttt{ca}: Neakceptované
   \item \texttt{bac}: Neakceptované
    \item \texttt{} (prázdny reťazec): Akceptované
\end{enumerate}




\newpage
\begin{itemize}
    \item \texttt{abcdef}: Tento regulárny výraz znamená, že symboly musia prichadzat´ postupne. Prazdny retazec nebude akceptovany.
\end{itemize}

\begin{verbatim}
 regular
   - sequence
   | -element
   |     -SYMBOL
   |         - a
   | -element
   |     -SYMBOL
   |         - b
   | -element
   |     -SYMBOL
   |         - c
   | -element
   |     -SYMBOL
   |         - d
   | -element
   |     -SYMBOL
   |         - e
   | -element
   |     -SYMBOL
   |         - f

   
    \end{verbatim}

Pre spracovanie regulárneho výrazu \texttt{abcdef} môžeme generovať konečný automat (NKA), ktorý akceptuje nasledovné reťazce:

\begin{enumerate}
    \item \texttt{ab}: Nekceptované
    \item \texttt{c}: Nekceptované
    \item \texttt{a}: Nekceptované
    \item \texttt{abcdef}: Akceptované
    \item \texttt{} (prázdny reťazec): Nekceptované
\end{enumerate}

\subsection{Problémy a ich riešenie}

% Opis problémov a riešení počas implementácie

\subsection{Hodnotenie zadania}

% Čo sa páčilo/nepáčilo

\section{Vyhodnotenie a záver}

% Zhrnutie a možné rozšírenia

\end{document}